\section{Vektorräume}

\subsection{Reeller Vektorraum}
Ein \emph{reeller Vektorraum} ist eine Menge $V \neq \emptyset$ mit Addition $+: V \times V \to V$ und skalarer Multiplikation $\cdot: \mathbb{R} \times V \to V$, für die gilt ($\vec{a}, \vec{b}, \vec{c} \in V$, $\lambda, \mu \in \mathbb{R}$):

\begin{enumerate}
    \item $\vec{a} + \vec{b} = \vec{b} + \vec{a}$ \quad (Kommutativ)
    \item $\vec{a} + (\vec{b} + \vec{c}) = (\vec{a} + \vec{b}) + \vec{c}$ \quad (Assoziativ)
    \item $\exists\, \vec{0} \in V: \vec{a} + \vec{0} = \vec{a}$ \quad (Neutralelement)
    \item $\forall\, \vec{a}\, \exists\, {-\vec{a}}: \vec{a} + (-\vec{a}) = \vec{0}$ \quad (Inverses)
    \item $\lambda(\mu \vec{a}) = (\lambda\mu)\vec{a}$ \quad (Assoziativ)
    \item $\lambda(\vec{a} + \vec{b}) = \lambda\vec{a} + \lambda\vec{b}$ \quad (Distributiv 1)
    \item $(\lambda + \mu)\vec{a} = \lambda\vec{a} + \mu\vec{a}$ \quad (Distributiv 2)
    \item $1 \cdot \vec{a} = \vec{a}$ \quad (Einselement)
\end{enumerate}

\textbf{Beispiele für Vektorräume:}
\begin{itemize}
    \item $\mathbb{R}^n$ (Vektoren mit $n$ Komponenten)
    \item $\mathbb{P}_n[x]$: Polynome vom Grad $\leq n$
    \item $\mathbb{R}^{m \times n}$: Menge aller $m \times n$-Matrizen
    \item $\mathbb{Z}^n/2$: Datenworte mit $n$ Bits (über Körper $\mathbb{Z}/2$)
\end{itemize}

\textbf{Kein} Vektorraum: Polynome vom Grad $= n$ (Addition kann Grad reduzieren).

\textbf{Abgeleitete Regeln:}
$0 \cdot \vec{a} = \vec{0}$, \quad $\lambda \cdot \vec{0} = \vec{0}$, \quad $(-1) \cdot \vec{a} = -\vec{a}$

\subsection{Unterräume}
Eine Teilmenge $U \subseteq V$ ist ein \emph{Unterraum} von $V$, wenn $U$ selbst ein Vektorraum ist.

\textbf{Unterraumkriterien} ($U \neq \emptyset$):
\begin{enumerate}
    \item $\vec{a}, \vec{b} \in U \Rightarrow \vec{a} + \vec{b} \in U$ \quad (abgeschlossen unter Addition)
    \item $\lambda \in \mathbb{R},\, \vec{a} \in U \Rightarrow \lambda\vec{a} \in U$ \quad (abgeschlossen unter Skalarmultiplikation)
\end{enumerate}

\textbf{Wichtig:} Falls $\vec{0} \notin U$, dann ist $U$ kein Unterraum.

\textbf{Beispiele:}
\begin{itemize}
    \item $\{\vec{0}\}$ ist immer ein Unterraum (\emph{Nullvektorraum})
    \item Geraden/Ebenen durch den Ursprung sind Unterräume von $\mathbb{R}^2$/$\mathbb{R}^3$
    \item Geraden/Ebenen \textbf{nicht} durch den Ursprung: \textbf{keine} Unterräume
    \item Symmetrische Matrizen $S^{m \times m} \subseteq \mathbb{R}^{m \times m}$: Unterraum
    \item $\mathbb{P}_2[x] \subseteq \mathbb{P}_4[x]$: Unterraum
    \item Lösungsmenge von $A\vec{x} = \vec{0}$: Unterraum von $\mathbb{R}^n$ (\emph{Kern} von $A$)
    \item Lösungsmenge von $A\vec{x} = \vec{c}$ mit $\vec{c} \neq \vec{0}$: \textbf{kein} Unterraum
\end{itemize}

\subsection{Linearer Spann}
Der lineare Spann ist die Menge aller Linearkombinationen und bildet einen Unterraum von $V$.

\[
    \operatorname{span}(\vec{b}_1, \ldots, \vec{b}_n) = \left\{ \sum_{k=1}^{n} \lambda_k \vec{b}_k \;\middle|\; \lambda_k \in \mathbb{R} \right\}
\]

\subsection{Erzeugendensystem}
$\{\vec{b}_1, \ldots, \vec{b}_n\}$ ist \emph{Erzeugendensystem} von $V$, wenn $V = \operatorname{span}(\vec{b}_1, \ldots, \vec{b}_n)$.

\textbf{Kriterium für $\mathbb{R}^m$:} Vektoren $\vec{b}_1, \ldots, \vec{b}_n \in \mathbb{R}^m$ als Spalten in $m \times n$-Matrix $B$ schreiben. Äquivalent:
\begin{enumerate}
    \item $\{\vec{b}_1, \ldots, \vec{b}_n\}$ ist Erzeugendensystem von $\mathbb{R}^m$
    \item $B \cdot \vec{x} = \vec{a}$ ist für jedes $\vec{a} \in \mathbb{R}^m$ lösbar
    \item $\operatorname{rg}(B) = m$
\end{enumerate}

\subsection{Lineare Unabhängigkeit}
$\vec{b}_1, \ldots, \vec{b}_n$ sind \emph{linear unabhängig}, wenn aus $\lambda_1 \vec{b}_1 + \ldots + \lambda_n \vec{b}_n = \vec{0}$ folgt: $\lambda_1 = \ldots = \lambda_n = 0$.

\textbf{Kriterium für $\mathbb{R}^m$:} Vektoren als Spalten in $m \times n$-Matrix $B$. Äquivalent:
\begin{enumerate}
    \item $\vec{b}_1, \ldots, \vec{b}_n$ sind linear unabhängig
    \item $B \cdot \vec{x} = \vec{0}$ hat nur die Lösung $\vec{x} = \vec{0}$
    \item $\operatorname{rg}(B) = n$
\end{enumerate}

\textbf{Zusammenhang:} Sind $\vec{b}_1, \ldots, \vec{b}_n$ linear unabhängig, so ist jede Linearkombination $\vec{a} = \lambda_1 \vec{b}_1 + \ldots + \lambda_n \vec{b}_n$ \textbf{eindeutig} (Koeffizienten eindeutig bestimmt).

\subsection{Basis und Dimension}
$\mathcal{B} = \{\vec{b}_1, \ldots, \vec{b}_n\}$ ist \emph{Basis} von $V$, wenn:
\begin{enumerate}
    \item $\mathcal{B}$ ist Erzeugendensystem von $V$
    \item $\vec{b}_1, \ldots, \vec{b}_n$ sind linear unabhängig
\end{enumerate}

\textbf{Satz für $\mathbb{R}^n$:} Vektoren $\vec{b}_1, \ldots, \vec{b}_n \in \mathbb{R}^n$ als Spalten in $n \times n$-Matrix $B$. Äquivalent:
\begin{enumerate}
    \item $\vec{b}_1, \ldots, \vec{b}_n$ bilden eine Basis von $\mathbb{R}^n$
    \item $\operatorname{rg}(B) = n$
    \item $\det(B) \neq 0$
    \item $B$ ist invertierbar
    \item $B \cdot \vec{x} = \vec{c}$ hat eine eindeutige Lösung
\end{enumerate}

Eine Basis von $\mathbb{R}^n$ besteht aus genau $n$ Vektoren. Für Erzeugendensystem braucht man $\operatorname{rg}(B) = m$, für lineare Unabhängigkeit $\operatorname{rg}(B) = n$ $\Rightarrow$ Basis erfordert $m = n$.

\textbf{Wichtige Basen:}
\begin{itemize}
    \item \emph{Standardbasis} von $\mathbb{R}^n$: $\mathcal{S} = \{\vec{e}_1, \ldots, \vec{e}_n\}$
    \item \emph{Monombasis} von $\mathbb{P}_n[x]$: $\mathcal{M} = \{1, x, x^2, \ldots, x^n\}$
\end{itemize}

\textbf{Dimension:} $\dim(V)$ = Anzahl Vekotren die eine Basis $V$ bilden.
\begin{itemize}
    \item $\dim(\mathbb{R}^n) = n$, \quad $\dim(\mathbb{P}_n[x]) = n + 1$
    \item $\dim(\mathbb{R}^{m \times n}) = m \cdot n$, \quad $\dim(S^{n \times n}) = \frac{n(n+1)}{2}$
    \item Unterraum $U \subset V$: $\dim(U) \leq \dim(V)$
    \item $\dim(\{\vec{0}\}) = 0$
\end{itemize}

\subsection{Komponentendarstellung und Basiswechsel}
Bezüglich Basis $\mathcal{B} = \{\vec{b}_1, \ldots, \vec{b}_n\}$: \quad $\vec{a} = \alpha_1 \vec{b}_1 + \ldots + \alpha_n \vec{b}_n \;\Rightarrow\; \vec{a} = \begin{pmatrix} \alpha_1 \\ \vdots \\ \alpha_n \end{pmatrix}_{\!\mathcal{B}}$

\textbf{$\mathcal{B} \to \mathcal{S}$:} Koeffizienten $\alpha_k$ und Basisvektoren $\vec{b}_k$ (in $\mathcal{S}$-Darstellung) einsetzen:
\[
    \vec{a}_{\mathcal{S}} = \alpha_1 \cdot \vec{b}_1 + \alpha_2 \cdot \vec{b}_2 + \ldots + \alpha_n \cdot \vec{b}_n = B \cdot \vec{a}_{\mathcal{B}}
\]

\textbf{$\mathcal{S} \to \mathcal{B}$:} LGS lösen: $B \cdot \vec{a}_{\mathcal{B}} = \vec{a}_{\mathcal{S}}$, wobei $B = (\vec{b}_1 \mid \vec{b}_2 \mid \ldots \mid \vec{b}_n)$.

